\begin{notechapter}
\chapter{Inner Products and Euclidean Affine Spaces}
\label{ch:inner-products}
\chaptersubtitle{%
  Riesz, the adjoint, and how a metric turns covectors into vectors}

\begin{chapterplan}
\begin{itemize}[nosep]
  \item Inner products, norms, Cauchy--Schwarz, orthogonality, Gram--Schmidt;
        the Gram matrix $M$ of an inner product in a basis.
  \item The Riesz representation theorem and the Riesz map $R_V:V\to V^*$.
  \item The adjoint $\adj A = R_V^{-1}\circ\dual A\circ R_W$: it maps vectors to
        vectors and depends on both inner products. Its matrix is $M^{-1}A^\top
        N$; it reduces to $A^\top$ exactly in orthonormal bases (conjugate
        transpose over $\C$).
  \item Dual versus adjoint, side by side: what each one needs, where each one
        lives, and why identifying $(\K^n)^*$ with $\K^n$ through the dual basis
        agrees with the Riesz identification only for an orthonormal basis.
  \item Orthogonal projections, least squares and the normal equations.
  \item Euclidean affine spaces: distances, angles and isometries.
\end{itemize}
\smallskip\noindent\textbf{Machine learning threads.}
\begin{itemize}[nosep]
  \item Preconditioning is a change of inner product.
  \item The gradient for the metric $M$ is $M^{-1}$ times the vector of partial
        derivatives (\cref{ch:gradient-hessian}).
  \item Whitening and the Mahalanobis distance.
\end{itemize}
\end{chapterplan}
\end{notechapter}
